\section{Definitions and Examples}

\begin{enumerate}
\item An \textbf{$\R$-linear equation} in the variables $x_1, \dots,
  x_n$ is an equation that can be expressed in the form
  \begin{displaymath}
    a_1 x_1 + \dots + a_n x_n = b
  \end{displaymath}
  where $a_1, \dots, a_n, b \in \R$.\footnote{In this course we will
  often drop the ``$\R$'' and simply call this a linear equation.}
  The numbers $a_1, \dots, a_n$ are called \textbf{coefficients} and
  the variables $x_1, \dots, x_n$ are also called \textbf{unknowns}.

  \textit{Examples.}
  \begin{itemize}
  \item $2x - 3y + \pi z = -44.5$ is a linear equation in the
    variables $x, y, z$.\footnote{The choice of variables names is
    unimportant, and will be clear from context.}
  \item $2(x_1 + x_2) = 4x_3 - x_4$ is a linear equation in the
    variables $x_1, x_2, x_3, x_4$.\footnote{This equation can be
    \textit{expressed} in the above form by rearranging terms.}
  \item $(2x_1 + x_2)(x_3 - x_4) = 5$ is not a linear equation because
    the left-hand-side, after rearranging terms, includes
    nonlinear terms, e.g., $x_2x_3$.
  \end{itemize}
\item A \textbf{system of linear equations} in the variables $x_1,
  \dots, x_n$ is a collection of linear equations in those variables.
  We also call this a \textbf{linear system}.  Note that every
  equation in a linear system must over the \textit{same} variables.
  The general form of a system of linear equations is written as
  follows.
  \begin{equation}
    \begin{split}
      a_{11}x_1 + \dots + a_{1n} x_n &= b_1 \\
      &\vdots \\
      a_{m1}x_1 + \dots + a_{mn} x_n &= b_m
    \end{split}
    \label{gen-form-lin-sys}
  \end{equation}

  \textit{Examples.}
  \begin{itemize}
  \item The following is a system of two linear equations over the
    variables $x$ and $y$.  This system represents a pair of lines in
    the plane that intersect at a single point.
    \begin{equation}
      \begin{split}
        2x + 3y &= 4 \\
        -5x + 3y &= 0
      \end{split}
      \label{lin-sys-ex-1}
    \end{equation}
  \item The following is a system of three linear equations over four
    variables.  Note that equations are required to have nonzero
    coefficients for every variable.
    \begin{equation}
      \begin{split}
      3x_1 - 4.4x_2 &= 1 \\
      200x_1 + 200x_2 + x_3 + 200x_4 &= 200 \\
      9x_1 - x_2 - 5x_4 &= 98
      \end{split}
      \label{lin-sys-ex-2}
    \end{equation}
  \end{itemize}
\item
  We write $\R^n$ for \textbf{$n$-dimensional Euclidean space}.  In
  this chapter we write $(a_1, \dots, a_n)$ where $a_1, \dots, a_n \in
  R$ for an element of $\R^n$.  We call this a \textbf{point} in
  $\R^n$.\footnote{In the next chapter, we will identify points in
  $\R^n$ with \textit{vectors}.}
\item
  An \textbf{$\R$-matrix} is a rectangular grid of real numbers.  We
  usually assume that a matrix has at least one row and and one
  column.  We write $\R^{m \times n}$ for the set of all matrices with
  $m$ rows and $n$ columns.  We identify $\R^{1 \by 1}$ with $\R$
  (i.e., we don't make a distinction between real numbers and matrices
  with one entry).

  \textit{Examples.}
  \begin{itemize}
  \item The following is a matrix in $\R^{3 \by 2}$.
    \begin{displaymath}
      \begin{bmatrix}
        1 & 2 \\
        \pi & -4 \\
        -22.4 & 0
      \end{bmatrix}
    \end{displaymath}
  \item The following is a matrix in $\R^{1 \by 5}$.
    \begin{displaymath}
      \begin{bmatrix}
        1 & 2 & 3 & 4 & 5
      \end{bmatrix}
    \end{displaymath}
  \end{itemize}
\item The \textbf{augmented matrix} of a system of $m$ linear
  equations over $n$ variables is the matrix in $\R^{m \by (n + 1)}$
  whose entries include the coefficients and right-hand-sides of each
  linear equation in the system.  The augmented matrix of the a
  general linear system (as in \ref{gen-form-lin-sys}) is as follows.
  \begin{displaymath}
    \begin{bmatrix}
      a_{11} & \dots & a_{1n} & b_1 \\
      \vdots & \ddots & \vdots & \vdots \\
      a_{m1} & \dots & a_{mn} & b_m
    \end{bmatrix}
  \end{displaymath}

  \textit{Examples.}
  \begin{itemize}
  \item
    The following matrix in $\R^{2 \by 3}$ is the augmented matrix of
    the linear system in \ref{lin-sys-ex-1}.
    \begin{displaymath}
      \begin{bmatrix}
        2 & 3 & 4 \\
        -5 & 3 & 0
      \end{bmatrix}
    \end{displaymath}
  \item
    The following matrix in $\R^{3 \by 5}$ is the augmented matrix of
    the linear system in \ref{lin-sys-ex-2}.
    \begin{displaymath}
      \begin{bmatrix}
        3 & -4.4 & 0 & 0 & 1 \\
        200 & 200 & 1 & 200 & 200 \\
        9 & -1 & 0 & -5 & 98
      \end{bmatrix}
    \end{displaymath}
  \end{itemize}
\item The \textbf{coefficient matrix} of a system of $m$ linear
  equations over $n$ variables is the matrix in $\R^{m \by n}$ gotten
  by removing the rightmost column of its augmented matrix.  The
  coefficient matrix of a general linear system (as in
  \ref{gen-form-lin-sys}) is as follows.
  \begin{align*}
    \begin{bmatrix}
      a_{11} & \dots & a_{1n} \\
      \vdots & \ddots & \vdots \\
      a_{m1} & \dots & a_{mn}
    \end{bmatrix}
  \end{align*}

  \textit{Examples.}
  \begin{itemize}
  \item
    The following matrix in $\R^{2 \by 2}$ is the coefficient matrix of
    the linear system in \ref{lin-sys-ex-1}.
    \begin{displaymath}
      \begin{bmatrix}
        2 & 3 \\
        -5 & 3
      \end{bmatrix}
    \end{displaymath}
  \item
    The following matrix in $\R^{3 \by 5}$ is the coefficient matrix of
    the linear system in \ref{lin-sys-ex-2}.
    \begin{displaymath}
      \begin{bmatrix}
        3 & -4.4 & 0 & 0 \\
        200 & 200 & 1 & 200 \\
        9 & -1 & 0 & -5
      \end{bmatrix}
    \end{displaymath}
  \end{itemize}
\item A \textbf{solution} of a linear system over the variables $x_1,
  \dots, x_n$ is a point $(s_1, \dots, s_n) \in \R^n$ such that
  substituting each variable $x_i$ with the value $s_i$ simultaneously
  satisfies every equation in the system.  A solution of the linear
  general system in \ref{gen-form-lin-sys} is a point $(s_1, \dots, s_n) \in
  \R^n$ such that the following real number equalities hold.
  \begin{align*}
    a_{11}s_1 + \dots + a_{1n} s_n &= b_1 \\
    &\vdots \\
    a_{m1}s_1 + \dots + a_{mn} s_n &= b_m \\
  \end{align*}

  \textit{Examples.}
  \begin{itemize}
  \item The point $\left(\frac{12}{21}, \frac{20}{21}\right)$ is a
    solution to the linear system in \ref{lin-sys-ex-1} because
    \begin{displaymath}
      2\left(\frac{12}{21}\right) + 3\left(\frac{20}{21}\right)
      = \frac{24}{21} + \frac{60}{21}
      = \frac{84}{21}
      = 4
      \quad
      \text{and}
      \quad
      -5\left(\frac{12}{21}\right) + 3\left(\frac{20}{21}\right)
      = -\frac{60}{21} + \frac{60}{21}
      = 0
    \end{displaymath}
  \item
    \todo{Another example.}
  \end{itemize}
\item A linear system is \textbf{consistent} if it has a solution.
  Otherwise, it is \textbf{inconsistent}.

  \textit{Examples.}
  \begin{itemize}
  \item As noted above, the linear system in \ref{lin-sys-ex-1} is consistent.
  \item The following linear system is inconsistent.  It represents
    two parallel lines in $\R^2$.
    \begin{align*}
      2x - 3y &= 4 \\
      -4x + 6y &= 5
    \end{align*}
  \end{itemize}

\item There are three kinds of \textbf{elementary row operations}.
  \begin{enumerate}
  \item \textbf{Replacement} operations add a scalar multiple of one
    row to another row.  They are denoted
    \begin{displaymath}
      R_i \gets R_i + cR_j
    \end{displaymath}
    where $c \in \R$ and act on matrices as follows.
    \begin{displaymath}
      \begin{bmatrix}
        a_{11} & \dots & a_{1n} \\
        \vdots & \ddots & \vdots \\
        a_{i1} & \dots & a_{in} \\
        \vdots & \ddots & \vdots \\
        a_{j1} & \dots & a_{jn} \\
        \vdots & \ddots & \vdots \\
        a_{m1} & \dots & a_{mn}
      \end{bmatrix}
      \xrightarrow{R_i \gets R_i + cR_j}
      \begin{bmatrix}
        a_{11} & \dots & a_{1n} \\
        \vdots & \ddots & \vdots \\
        a_{i1} + ca_{j1} & \dots & a_{in} + ca_{jn} \\
        \vdots & \ddots & \vdots \\
        a_{j1} & \dots & a_{jn} \\
        \vdots & \ddots & \vdots \\
        a_{m1} & \dots & a_{mn}
      \end{bmatrix}
    \end{displaymath}
  \item \textbf{Scaling} operations multiply a row by a nonzero
    constant. They are denoted
    \begin{displaymath}
      R_i \gets cR_i
    \end{displaymath}
    where $c \in \R$ and $c \neq 0$ and act on matrices as follows.
    \begin{displaymath}
      \begin{bmatrix}
        a_{11} & \dots & a_{1n} \\
        \vdots & \ddots & \vdots \\
        a_{i1} & \dots & a_{in} \\
        \vdots & \ddots & \vdots \\
        a_{m1} & \dots & a_{mn}
      \end{bmatrix}
      \xrightarrow{R_i \gets cR_i}
      \begin{bmatrix}
        a_{11} & \dots & a_{1n} \\
        \vdots & \ddots & \vdots \\
        ca_{i1} & \dots & ca_{in} \\
        \vdots & \ddots & \vdots \\
        a_{m1} & \dots & a_{mn}
      \end{bmatrix}
    \end{displaymath}
  \item \textbf{Swap} operations interchange two rows.  They are
    denoted
    \begin{displaymath}
      R_i \leftrightarrow R_j
    \end{displaymath}
    and act on matrices as follows.
    \begin{displaymath}
      \begin{bmatrix}
        a_{11} & \dots & a_{1n} \\
        \vdots & \ddots & \vdots \\
        a_{i1} & \dots & a_{in} \\
        \vdots & \ddots & \vdots \\
        a_{j1} & \dots & a_{jn} \\
        \vdots & \ddots & \vdots \\
        a_{m1} & \dots & a_{mn}
      \end{bmatrix}
      \xrightarrow{R_i \leftrightarrow R_j}
      \begin{bmatrix}
        a_{11} & \dots & a_{1n} \\
        \vdots & \ddots & \vdots \\
        a_{j1} & \dots & a_{jn} \\
        \vdots & \ddots & \vdots \\
        a_{i1} & \dots & a_{in} \\
        \vdots & \ddots & \vdots \\
        a_{m1} & \dots & a_{mn}
      \end{bmatrix}
    \end{displaymath}
  \end{enumerate}

  \textit{Examples.}
  \begin{itemize}
  \item The following is an example of applying the replacement
    operation $R_2 \gets R_2 - 3R_1$.
    \begin{displaymath}
      \begin{bmatrix}
        1 & 2 & 3 \\
        4 & 5 & 6 \\
        7 & 8 & 9
      \end{bmatrix}
      \xrightarrow{R_2 \gets R_2 - 3R_1}
      \begin{bmatrix}
        1 & 2 & 3 \\
        1 & -1 & -3 \\
        7 & 8 & 9
      \end{bmatrix}
    \end{displaymath}
  \item The following is an example of applying the scaling operation $R_3 \gets (-4)R_3$.
    \begin{displaymath}
      \begin{bmatrix}
        1 & 2 & 3 \\
        4 & 5 & 6 \\
        7 & 8 & 9
      \end{bmatrix}
      \xrightarrow{R_3 \gets (-4)R_3}
      \begin{bmatrix}
        1 & 2 & 3 \\
        4 & 5 & 6 \\
        -28 & -32 & -36
      \end{bmatrix}
    \end{displaymath}
  \item The following is an example of applying the swap operation $R_1 \leftrightarrow R_2$.
    \begin{displaymath}
      \begin{bmatrix}
        1 & 2 & 3 \\
        4 & 5 & 6 \\
        7 & 8 & 9
      \end{bmatrix}
      \xrightarrow{R_1 \leftrightarrow R_2}
      \begin{bmatrix}
        4 & 5 & 6 \\
        1 & 2 & 3 \\
        7 & 8 & 9
      \end{bmatrix}
    \end{displaymath}
  \end{itemize}
\item Matrices $A, B \in \R^{m \times n}$ are \textbf{row equivalent},
  written $A \sim B$, if there is a sequence of elementary row
  operations that transform $A$ into $B$.

  \textit{Facts.}
  \begin{itemize}
  \item
    \textit{Row equivalence is, in fact, an equivalence.}
    \begin{itemize}
    \item (reflexivity) $A \sim A$ for any matrix $A$;
    \item (symmetry) $A \sim B$ if and only if $B \sim A$;
    \item (transitivity) If $A \sim B$ and $B \sim C$ then $A \sim C$.
    \end{itemize}
  \item
    \textit{Row operations preserve solutions.}  Suppose $A$ and $B$
    are the augmented matrices of two linear systems.  If $A \sim B$
    then $\vs \in \R^n$ is a solution of the linear system with
    augmented matrix $A$ if and only if it is a solution of the system
    with augmented matrix $B$.
  \end{itemize}

  \textit{Example.}
  \todo{Fill in example.}
\item
  The \textbf{leading entry} of a row in a matrix is its first nonzero
  entry.  Note that a row may not have a leading entry.

  \textit{Example.}
  \todo{Fill in example.}
\item
  A matrix is in \textbf{echelon form} if
  \begin{itemize}
  \item
    the leading entry of every row appears to the right of the
    leading entries of the rows that precede it;
  \item
    every all-zero row appears below every non-zero row.
  \end{itemize}
  This can be expressed by example pictorially as
  \begin{displaymath}
    \begin{bmatrix}
      0&\blacksquare&*&*&*&*&*&*&*&*\\
      0&0&0&\blacksquare&*&*&*&*&*&*\\
      0&0&0&0&\blacksquare&*&*&*&*&*\\
      0&0&0&0&0&\blacksquare&*&*&*&*\\
      0&0&0&0&0&0&0&0&\blacksquare&*\\
      0&0&0&0&0&0&0&0&0&0\\
    \end{bmatrix}
  \end{displaymath}
  where ``$\blacksquare$'' represents a nonzero entry and ``$*$" represents
  an entry of any value.

  \textit{Facts.}
  \begin{itemize}
  \item
    \textit{Echelon forms can be used to determine consistency.}
    Let $A$ be the augmented matrix of a linear system.  The system is
    inconsistent if and only if every echelon form of $A$ has a row of
    the form
    \begin{displaymath}
      \begin{bmatrix}
        0 \dots & 0 & b
      \end{bmatrix}
    \end{displaymath}
    where $b \neq 0$.
  \end{itemize}

  \textit{Examples.}
  \begin{itemize}
  \item
    \todo{Example}
    \todo{Example}
  \end{itemize}
\item A matrix is in \textbf{row-reduced echelon form} (RREF) if it is echelon form and
  \begin{itemize}
  \item every leading entry is $1$;
  \item its leading entries are the only nonzero entries of their columns.
  \end{itemize}
  This can be expressed by example pictorially as
  \begin{displaymath}
    \begin{bmatrix}
      0&1&*&0&0&0&*&*&0&*\\
      0&0&0&1&0&0&*&*&0&*\\
      0&0&0&0&1&0&*&*&0&*\\
      0&0&0&0&0&1&*&*&0&*\\
      0&0&0&0&0&0&0&0&1&*\\
      0&0&0&0&0&0&0&0&0&0\\
    \end{bmatrix}
  \end{displaymath}
  where ``$*$" represents an entry of any value.  We will also say
  that a matrix \textit{is} an RREF to mean that it is in row-reduced
  echelon form.

  \textit{Fact.} \textit{Every matrix is row-equivalent to exactly one RREF.}

  \textit{Examples.}
  \begin{itemize}
  \item \todo{Example}
  \item \todo{Example}
  \end{itemize}

\item A \textbf{pivot position} of a matrix is a position of a leading
  entry in its RREF.  A \textbf{pivot column} is a column containing a
  pivot position.
\item A variable $x_i$ is \textbf{basic} with respect to a linear
  system if the $i$th column of its augmented matrix is a pivot
  column.  A variables is \textbf{free} if it is not basic.
\item A \textbf{general form solution}\footnote{This is also
sometimes just called a general solution.} is a collection of statements of the form
  \begin{displaymath}
    x_i = x_i \text{ is free} \qquad \text{or} \qquad x_i = a_1x_{j_1} + \dots + a_kx_{j_k}
  \end{displaymath}
  where $a_{1}, \dots, a_{k} \in \R$ and $x_{j_1}, \dots, x_{j_k}$ are
  free variables.  In other words, a general form solution is one that
  expresses every basic variable of a linear system linearly in terms
  of its free variables.
\end{enumerate}

\subsection*{Topics Not Covered}

\begin{itemize}
\item Gaussian elimination
\item deriving general form solutions from RREFs
\end{itemize}


\section{Basic Exercises}

Determine the coefficient matrix and augmented matrix for each of the
following linear systems.
\begin{tasks}[style=enumerate](2)
\task
  \begin{align*}
    2 x_{1} - 6 x_{2} &= -4 \\
    6 x_{1} + 8 x_{2} &= -7
  \end{align*}
\task
  \begin{align*}
    x_{1} + 8 x_{2} + 7 x_{3} &= -1 \\
    4 x_{1} - 9 x_{2} &= 8 \\
    7 x_{1} - 7 x_{2} - 3 x_{3} &= 10
  \end{align*}
\task
  \begin{align*}
    x_{1} - 2x_{2} - 2x_{3} &= 2 \\
    2x_{1} - 3x_{2} - 5x_{3} &= 2 \\
    -2x_{1} + 2x_{2} + 7x_{3} &= -1
  \end{align*}
\task
  \begin{align*}
    - 8 x_{1} + 6 x_{2} + x_{3} + 7 x_{4} + 10 x_{5} &= 8 \\
    - 4 x_{2} - 4 x_{3} - x_{4} + x_{5} &= -8 \\
    7 x_{1} + x_{2} - 2 x_{3} - 8 x_{4} + 7 x_{5} &= 9 \\
    - 4 x_{1} - 7 x_{2} + 5 x_{3} + 2 x_{4} - 9 x_{5} &= -4
  \end{align*}
\end{tasks}
For each of the following linear systems, verify that the given point
$s$ is a solution.
\begin{tasks}[resume, style=enumerate](2)
\task
  \begin{align*}
    x_{1} - 2x_{2} + x_{3} - 2x_{4} &= -9 \\
    x_{1} - x_{2} - x_{3} - 2x_{4} &= -10 \\
    -3x_{1} + 8x_{2} - 6x_{3} + 4x_{4} &= 21 \\
    2x_{2} - 7x_{3} + 7x_{4} &= 13 \\
    \\
    s = (1, 3, 2, 3)
  \end{align*}
\task
  \begin{align*}
    2 x_{1} - 2 x_{2} - 6 x_{3} + 4 x_{4} &= -20 \\
    x_{1} + 2 x_{2} - 6 x_{3} - 4 x_{4} &= 26 \\
    x_{1} - 4 x_{3} - x_{4} &= 5 \\
    \\
    s = (10, 8, 2, -3)
  \end{align*}
\end{tasks}
Demonstrate that each of the following linear systems has a unique
solution.
\begin{tasks}[resume, style=enumerate](3)
\task
  \begin{align*}
    3 x_{1} + 6 x_{2} &= -15 \\
    - x_{1} + x_{2} &= 5
  \end{align*}
\task
  \begin{align*}
    2x + 4y &= 29 \\
    -6x - 10y &= -75
  \end{align*}
\task
  \begin{align*}
    x + 3y - 3z &= -2 \\
    2x + 7y + 7z &= -5 \\
    x - 2y - z &= 3
  \end{align*}
\task
  \begin{align*}
    3x - 4y + 3z &= -9 \\
    6x + 7y - 3z &= 0 \\
    x + 10z &= -21
\end{align*}
\task
  \begin{align*}
    - 3 x_{1} + 6 x_{3} &= 12 \\
    3 x_{1} + 3 x_{2} - 12 x_{3} &= -21 \\
    2 x_{1} - 6 x_{2} + 10 x_{3} &= 16
  \end{align*}
\task
  \begin{align*}
    - 3 x_{1} + 6 x_{2} &= 9 \\
    3 x_{1} - 8 x_{2} &= -9 \\
    2 x_{2} + 6 x_{3} &= 18
  \end{align*}
\task
  \begin{align*}
    x_{1} - 2x_{2} - 2x_{3} &= -7 \\
    -x_{1} + 3x_{2} + 2x_{3} &= 10 \\
    2x_{1} - 6x_{2} - 3x_{3} &= -18
  \end{align*}
\task
  \begin{align*}
    - 3 x_{1} - 6 x_{3} + 9 x_{4} &= -6 \\
    - x_{1} + 2 x_{2} - 3 x_{4} &= 8 \\
    - 2 x_{1} - 9 x_{2} - 19 x_{3} + 45 x_{4} &= -61 \\
    x_{1} - 3 x_{2} + 5 x_{3} - 9 x_{4} &= 5
  \end{align*}
\end{tasks}
For each of the following matrices, apply the given row operations
from top to bottom.
\begin{tasks}[resume, style=enumerate](2)
\task
  \begin{displaymath}
    \begin{bmatrix}
      -2 & 7 & 9 \\
      7 & -9 & -9 \\
      4 & -1 & 9
    \end{bmatrix}
  \end{displaymath}
  \begin{align*}
    R_{1} &\gets -5R_{1} \\
    R_{1} &\gets R_{1} - 4R_{2} \\
    R_{2} &\gets R_{2} + 3R_{1}
  \end{align*}
\task
  \begin{displaymath}
    \begin{bmatrix}
      -4 & 1 & 3 \\
      -7 & 0 & 5 \\
      -4 & 3 & 1
    \end{bmatrix}
  \end{displaymath}
  \begin{align*}
    R_{3} &\gets R_{3} - 3R_{2} \\
    R_{3} &\gets R_{3} + 3R_{1} \\
    R_{1} &\leftrightarrow R_{3}
  \end{align*}
\task
  \begin{displaymath}
    \begin{bmatrix}
      9 & 5 & -7 & -5 & -9 \\
      5 & -7 & 1 & -2 & -9 \\
      5 & 1 & -10 & 6 & -5 \\
      5 & 7 & -5 & 2 & 1
    \end{bmatrix}
  \end{displaymath}
  \begin{align*}
    R_{4} &\gets -R_{4} \\
    R_{2} &\gets R_{2} - 2R_{3} \\
    R_{2} &\gets R_{2} - 5R_{4} \\
    R_{3} &\gets R_{3} + 3R_{4} \\
    R_{3} &\leftrightarrow R_{2}
  \end{align*}
\end{tasks}
Determine a general form solution from each of the following
matrices. That is, determine a general form solution for a linear
system whose augmented matrix is row equivalent to the given matrix.
\begin{tasks}[resume, style=enumerate](2)
\task
  \begin{displaymath}
    \begin{bmatrix}
      1 & -3 & 0 & 0 & -5 \\
      0 & 0 & 1 & 0 & 8 \\
      0 & 0 & 0 & 1 & 5
    \end{bmatrix}
  \end{displaymath}
\task
  \begin{displaymath}
    \begin{bmatrix}
      1 & 0 & -3 & 0 & 1 \\
      0 & 1 & 1 & 1 & 1 \\
      0 & 0 & 0 & 1 & 5
    \end{bmatrix}
  \end{displaymath}
\task
  \begin{displaymath}
    \begin{bmatrix}
      1 & 0 & -6 & 0 & 3 \\
      0 & 1 & -3 & 2 & -9 \\
      0 & 0 & 0 & 1 & -6
    \end{bmatrix}
  \end{displaymath}
\task
  \begin{displaymath}
    \begin{bmatrix}
      1 & 1 & 0 & 0 & 0 & -4 & 5 & -3 \\
      0 & 0 & 1 & 0 & 0 & 1 & 3 & -4 \\
      0 & 0 & 0 & 0 & 1 & 5 & -3 & 2 \\
      0 & 0 & 0 & 0 & 0 & 0 & 0 & 0 \\
      0 & 0 & 0 & 0 & 0 & 0 & 0 & 0
    \end{bmatrix}
  \end{displaymath}
\end{tasks}
Determine three particular solutions of the linear system underlying
each of the following matrices.  That is, determine three particular
solutions of a linear system whose augmented matrix is row equivalent
to the given matrix.
\begin{tasks}[resume, style=enumerate](2)
\task
  \begin{displaymath}
    \left[\begin{matrix}1 & 2 & 0\\0 & 0 & 1\end{matrix}\right]
  \end{displaymath}
\task
  \begin{displaymath}
    \left[\begin{matrix}1 & 0 & 0 & -3 & -6\\0 & 1 & 0 & 5 & 5\\0 & 0 & 1 & 0 & -4\end{matrix}\right]
  \end{displaymath}
\task
  \begin{displaymath}
    \begin{bmatrix}
      1 & 1 & 0 & -3 & 0 & 0 & -3 & 1 \\
      0 & 0 & 1 & -1 & 0 & 3 & 2 & 3 \\
      0 & 0 & 0 & 0 & 1 & 3 & -4 & -3 \\
      0 & 0 & 0 & 0 & 0 & 0 & 0 & 0 \\
      0 & 0 & 0 & 0 & 0 & 0 & 0 & 0
    \end{bmatrix}
  \end{displaymath}
\end{tasks}
Determine the row-reduced echelon form of each of the following
matrices.  You must include the intermediate matrices and row
operations used.
\begin{tasks}[resume, style=enumerate](2)
\task
  \begin{displaymath}
    \begin{bmatrix}
      5 & 2 & 9 \\
      7 & 3 & 12 \\
      2 & 1 & 3
    \end{bmatrix}
  \end{displaymath}
\task
  \begin{displaymath}
    \begin{bmatrix}
      0 & 1 & -6 & -2 \\
      1 & 1 & -3 & -6 \\
      -2 & -4 & 18 & 16
    \end{bmatrix}
  \end{displaymath}
\task
  \begin{displaymath}
    \begin{bmatrix}
      2 & 3 & -13 & -8 \\
      1 & 1 & -4 & -1 \\
      1 & 0 & 1 & 5
    \end{bmatrix}
  \end{displaymath}
\task
  \begin{displaymath}
    \begin{bmatrix}
      1 & -1 & -2 & 1 \\
      -1 & 2 & 4 & 0 \\
      2 & -3 & -6 & 2 \\
      -2 & 1 & 2 & -1
    \end{bmatrix}
  \end{displaymath}
\end{tasks}
Determine a general form solution for each of the following linear
systems by first determining the row-reduced echelon form of its
augmented matrix.  If this system has no solutions then write
\enquote{no solution}.
\begin{tasks}[resume, style=enumerate](2)
\task
  \begin{align*}
    x - 2y - 7z &= -3 \\
    2x + y + 6z &= 14 \\
    x + y + 5z &= 9
  \end{align*}
\task
  \begin{align*}
    x_1 + x_2 - x_3 &= -9 \\
    x_2 - 2x_3 &= -1 \\
    x_1 + x_2 &= -10 \\
  \end{align*}
\task \faCalculator
  \begin{align*}
    x_1 - 2x_2 + 5x_3 &= 6 \\
    -2x_1 + 6x_2 - 11x_3 + 7x_4 &= -8 \\
    5x_1 - 10 x_2 + 25x_3 + 3x_4 &= 30
  \end{align*}
\task \faCalculator
  \begin{align*}
    x_1 - 3x_2 + 4x_3 &= 0 \\
    -x_1 + 6x_2 + x_3 &= 4 \\
    23x_2 + 5x_3 &= 9
  \end{align*}
\task \faCalculator
  \begin{align*}
    x_1 + 5x_2 + 3x_3 &= -4 \\
    x_1 + 6x_2 + 7x_3 &= -13 \\
    -2x_1 - 12 x_2 - 14 x_3 &= 25
  \end{align*}
\task
  \begin{align*}
    -x - 2y - z &= -4 \\
    x + 2y &= 3 \\
    -2x -4y + z &= -7
  \end{align*}
\task
  \begin{align*}
    x_{1} - 6x_{2} + x_{4} &= 2 \\
    2x_{1} - 12x_{2} + x_{3} - 4x_{4} &= 7 \\
    x_{1} - 6x_{2} - 2x_{3} + 13x_{4} &= -4
  \end{align*}
\task \faCalculator
  \begin{align*}
    88x_1 + 61x_2 + 8x_3 + 61 x_4 &= -68 \\
    37x_1 + 14x_2 + 72x_3 + 41 x_4 &= -460 \\
    57x_1 + 59x_2 + 69x_3 + 75x_4 &= -333 \\
    92 x_1 + 60x_2 + 28x_3 + 72x_4 &= -192
  \end{align*}
\task \faCalculator
  \begin{align*}
    x_2 + x_3 + x_4 + 3x_5 - x_6 + 2x_7 &= 34 \\
    x_3 - x_4 + 4x_5 &= -14 \\
    2x_2 - 4x_3 + 8x_4 - 18x_5 - x_6 &= 110 \\
    3x_3 - 3x_4 + 12 x_5 + 2x_7 &= -18 \\
    x_2 + 4x_3 - 2x_4 + 15x_5 - x_6 + 7x_7 &= 52
  \end{align*}
\task \faCalculator
  \begin{align*}
    8x_1 + 3x_2 - 3x_3 &= -4 \\
    -3x_1 + 8x_2 + 5x_3 &= -4 \\
    6x_1 + 7x_2 - 3x_3 &= -3 \\
    8x_1 + x_2 - x_3 &= -4
  \end{align*}
\end{tasks}
Determine a sequence of elementary row operations from the left matrix
to the right matrix.  You must include the intermediate matrices and
row operations used in your solution.
\begin{tasks}[resume, style=enumerate](2)
\task
  \begin{displaymath}
    \begin{bmatrix}
      -4 & -7 & 4 \\
      3 & 8 & 2 \\
      -10 & 1 & -9
    \end{bmatrix}
    \sim
    \begin{bmatrix}
      3 & 8 & 2 \\
      -14 & -6 & -5 \\
      -10 & 1 & -9
    \end{bmatrix}
  \end{displaymath}
\task
  \begin{displaymath}
    \begin{bmatrix}
      1 & 2 & -1 & 5 \\
      0 & 1 & 0 & 1 \\
      0 & 0 & 0 & 5 \\
      2 & 2 & 1 & 1
    \end{bmatrix}
    \sim
    \begin{bmatrix}
      3 & 4 & 0 & 6 \\
      4 & 4 & 2 & 2 \\
      0 & -1 & 0 & 4 \\
      0 & 1 & 0 & 1
    \end{bmatrix}
  \end{displaymath}
\end{tasks}

\section{True/False}

\trueFalseHeader

\begin{tasks}[resume, style=enumerate](1)
\task
  Elementary row operations cannot change the solution set of a linear
  system.
\task
  There is a linear system with exactly three solutions.
\task
  If $A$ is the augmented matrix of an inconsistent linear system, and
  $B$ is a matrix such that $A \sim B$ (that is, $A$ and $B$ are row
  equivalent), then $B$ is the augmented matrix of an inconsistent
  linear system.
\task
  If $A$ is the augmented matrix of an inconsistent system, then $A$
  has a pivot in its rightmost column.
\task
  If $A \sim B$ and $A \sim C$ and $B$ and $C$ are in reduced echelon
  form, then $B = C$.
\task
  There is a unique sequence of row operations that reduces a matrix
  to reduced echelon form.
\task
  If a general form solution of a linear system has a free variable,
  then the system has infinitely many solutions.
\task
  A matrix may have different pivot positions depending on the
  sequence of row operations used to attain a matrix in echelon form.
\task
  Every matrix has a unique echelon form.
\task
  A linear system over 3 variables and 2 equations must be consistent.
\task
  If the coefficient matrix of a linear system as more rows than
  columns, then the system has infinitely many solutions.
\task
  If $A$ is the augmented matrix of a linear system and it has a pivot
  position in every column, then the system is inconsistent.
\task
  A \textit{consistent} linear system whose augmented matrix is a $207
  \times 209$ matrix has infinitely many solutions.
\task
  A \textit{consistent} linear system whose coefficient matrix is
  square has infinitely many solutions.
\task
  If the rightmost column of the augmented matrix of a linear system
  is a pivot column, then the system is inconsistent.
\task
  If a system of linear equations of linear equations has infinitely
  many solutions, then it must have more unknowns than equations.
\task
  For any nonzero real values $a$ and $b$, the matrix
  $
  \begin{bmatrix}
    a & 2a \\
    b & 3b
  \end{bmatrix}
  $
  has a pivot in every column and every row.
\task
  A system of linear equations with 3 variables and 4 equations cannot
  be consistent.
\end{tasks}

\section{More Difficult Problems}
\begin{tasks}[resume, style=enumerate](1)
\task
  \faCalculator \ Consider the following linear system.
  \begin{align*}
    4x_2 + x_3 &= 16 \\
    9x_1 - 20x_2 - 8x_3 &= -71 \\
    3x_1 - 8x_2 - 3x_3 &= -29
  \end{align*}
  \begin{enumerate}
  \item Determine the RREF of the augmented matrix of the above linear
    system.
  \item Determine a general form solution for the above linear system.
  \item
    Find a solution to the linear system below with the following
    properties:
    \begin{enumerate}
    \item
      the solution consists entirely of integer values;
    \item
      the values in the solution are relatively prime, i.e., it's not
      possible to divide the solution by a number to get another
      integer solution. So $(2, 4, 6)$ does not satisfy this property
      but $(1, 2, 3)$ does.\footnote{This is the process we would use
      to find a valid solution to the problem of balancing chemical
      equations.}
    \end{enumerate}
  \end{enumerate}
\task
  Determine all the $2 \times 2$ matrices in echelon form whose
  entries are either $0$ or $1$.  Mark which ones are in reduced
  echelon form.
\task
  Determine every $3 \times 3$ matrix in reduced echelon form with at
  least two pivot positions whose entries are either $0$ or $1$.
\task
  Determine a linear system over three variables such that $(a, b, c)$
  is a solution exactly when the cubic function $f(x) = ax^3 + bx^2 +
  c$ intersects the points $(-1, 4)$, $(1, 5)$ and $(2, 10)$.  You do
  not need to solve the linear system.
\task
  Determine the slope-intercept form of the line equation which
  defines the intersection of the plane
  \begin{align*}
    2x + 3y + 3z = 6
  \end{align*}
  with the $xy$-plane.
\task
  For what values of the coefficient $h$ is the following system
  inconsistent?
  \begin{align*}
    x + 4y &= -1 \\
    3x - hy &= 7
  \end{align*}
  Is there a value of $h$ for which the above system has infinitely
  many solutions? Justify your answer.
\task
  Consider the following linear system with two unknown coefficients
  $h$ and $k$.
  \begin{align*}
    hx + 2y &= 1 \\
    3x + 9y &= k
  \end{align*}
  \begin{enumerate}
  \item
    Determine values of $h$ and $k$ so that the above linear system
    has no solutions.
  \item
    Determine values of $h$ and $k$ so that the above linear system
    has exactly one solution.
  \item
    Determine values of $h$ and $k$ so that the above linear system
    has infinitely many solutions.
  \end{enumerate}
\task
  Consider the following general form solution
  \begin{align*}
    x_1 &= 3x_2 - 2x_3 \\
    x_2 &\qquad \text{is free} \\
    x_3 &\qquad \text{is free} \\
    x_4 &= 2 + x_3
  \end{align*}
  \begin{enumerate}
  \item
    Determine another general form solution which describes the same
    solution set, but for which $x_3$ is basic and $x_4$ is free.
  \item
    Determine an RREF which yields the general form solution from the
    previous part.
  \end{enumerate}
\task
  Consider the following general form solution.
  \begin{align*}
    x_{1} &= -6 + 6x_{3} + 2x_{5} \\
    x_{2} &= 4 + 4x_{3} + 6x_{5} \\
    x_{3} &\text{ is free} \\
    x_{4} &= -4 + 5x_{5} \\
    x_{5} &\text{ is free}
  \end{align*}
  Determine a general form solution that describes the same solution
  set but in which $x_1$ is free.
\task
  \faCalculator \ Suppose you're investigating a claim that, out of
  five competing car companies who purchase engine parts from the same
  suppliers, one purportedly \textit{overestimated} their total
  spending by \textdollar 1,300,000.  Companies are required to report
  their total spending, and you've been able to determine how many
  units (× 10,000) of each part that each company has purchased.  You
  haven't been unable to determine how much each item costs per unit
  (it's an industry secret), but you can assume that each company pays
  the same amount per unit.  Given the following data, which company
  is falsifying their records?  Justify our answer.  The total amount
  spent by each company in the table below is multiplied by
  \textdollar 100,000.
  \begin{center}
    \begin{tabular}{|c|r|r|r|r|r|r|}
      \hline
      Co. & Part 1 & Part 2 & Part 3 & Part 4 & Part 5 & Total Spent\\
      \hline
      A & 15 & 3 & 24 & 46 & 182 & 1013 \\
      B & 5 & 3 & 14 & 25 &100 & 552 \\
      C & 15 & 3 & 24 & 46 &188 & 1038 \\
      D & 5 & 0 &  5 & 10 & 40 & 225 \\
      E & 15 & 3 & 26 & 47 &190 & 1056 \\
      \hline
    \end{tabular}
  \end{center}
\task
  Suppose you're given a system of linear equations with three
  variables and two equations, and that $(4, 1, 0)$ and $(2, 4, 1)$
  are solutions to this system.  You are also given that the two
  equations define distinct planes in $\R^3$.  Determine the
  RREF of the augmented matrix for this system.
\task
  Repeat the previous problem with the points $(-3, 0, 1)$ are $(-3,
  1, 1)$.
\task
  Consider an arbitrary system of linear equations with $n$ unknowns
  and $m$ equations.  Further suppose that
  \begin{itemize}
  \item
    it has a unique solution;
  \item
    it has at least as many equations as unknowns ($m \geq n$).
  \end{itemize}
  Write down an expression in terms of $m$ and $n$ for the number of
  all-zero rows which appear in the rwo-reduced echelon form of its
  augmented matrix.  Justify your answer.
\end{tasks}

\section{Challenge Problems}

\begin{tasks}[resume, style=enumerate](1)
\task
  Consider the following pair of linear systems.
  \begin{align*}
    ax + by &= c \\
    dx + ey &= f
  \end{align*}
  \begin{align*}
    ax + by - cz &= 0 \\
    dx + ey - fz &= 0
  \end{align*}
  \begin{enumerate}
  \item
    Demonstrate that if the first system has a solution, then so does
    the second one.
  \item
    Give \textbf{nonzero} values to $a$ through $f$ such that the
    second system has a solution, but the first does not.  Present
    your solution as an augmented matrix for the first system, i.e.,
    of the form
    \begin{displaymath}
      \begin{bmatrix}
        a & b & c \\
        d & e & f
      \end{bmatrix}
    \end{displaymath}
  \end{enumerate}
\task
  Consider the following pair of linear systems.
  \begin{align*}
    x_1 - 2x_2 &= 3 \\
    4x_1 + x_2 &= 21
  \end{align*}
  \begin{align*}
    10x_3 + 2x_4 &= x_1 \\
    -8x_3 + 9x_4 &= x_2
  \end{align*}
  \begin{enumerate}
  \item
    Solve the first system of linear equations (in $x_1$ and $x_2$)
    and write down the augmented matrix of the second system with the
    solutions of $x_1$ and $x_2$ substituted in.
  \item
    Determine the augmented matrix of a \textit{single} system of
    linear equations \textit{with all four equations} in the variables
    $x_1$, $x_2$, $x_3$, $x_4$.  Describe the relationship between
    this matrix and the one in the previous part.
\end{enumerate}
\task
  Determine what must hold of $a$, $b$, $c$, $d$ so that the following
  linear system has exactly one solution.
  \begin{align*}
    ax + by &= 0 \\
    cx + dy &= 0
  \end{align*}
\task
  Determine an inconsistent linear system in 3 variables such that
  every pair of equations is consistent.  That is, give values for $a$
  through $l$ in the system
  \begin{align*}
    ax + by + cz &= d \\
    ex + fy + gz &= h \\
    ix + jy + kz &= l
  \end{align*}
  such that the system is inconsistent, but each pair of equations
  forms a consistent system.  Present your solution to each part as
  the an augmented matrix, i.e., of the form
  \begin{displaymath}
    \begin{bmatrix}
      a & b & c & d \\
      e & f & g & h \\
      i & j & k & l
    \end{bmatrix}
  \end{displaymath}
  \begin{enumerate}
  \item
    Achieve this with no more than 5 nonzero values for $a$ through
    $l$.
  \item
    Achieve this with all nonzero values for $a$ through $l$.
  \end{enumerate}
\task
  Consider an arbitrary linear system in three variables:
  \begin{align*}
    ax + by + cz &= d \\
    ex + fy + gz &= h \\
    ix + jy + kz &= l
  \end{align*}
  Show that if $(s_x, s_y, s_z)$ and $(t_x, t_y, t_z)$ are solutions
  to the above system, then
  \begin{displaymath}
    \left(
    \frac{s_x + t_x}{2},
    \frac{s_y + t_y}{2},
    \frac{s_z + t_z}{2}
    \right)
  \end{displaymath}
  is also a solution.  Describe why $(s_x, s_y, s_z) \not = (t_x, t_y,
  t_z)$ implies the system has infinitely many solutions.
\task
  One way to describe a line in 3 dimensions is to use a parameter
  $t$.  That is, after fixing values $a_x$, $b_x$, $a_y$, $b_y$,
  $a_z$, and $b_z$, a line can be described as all points of the form
  \begin{displaymath}
    (a_xt + b_x, a_yt + b_y, a_zt + b_z)
  \end{displaymath}
  for any real value of $t$.  Give a pair of 3 dimensional linear
  equations (each of which represents a plane in $\R^3$) whose
  intersection is exactly the line whose points are defined by the
  above parametric form.  \textit{Hint:} The line above can be thought
  of as a system in the variables $x$, $y$, $z$, and $t$, e.g., one of
  its equations is $x - a_xt = b_x$. Write $t$ in terms of $x$ and
  substitute this value for $t$ into the other equations.
\task
  There are several ways to define row equivalence.  For example,
  suppose we restrict the replacement rule to only allow operations of
  the form
  \begin{displaymath}
    R_i \gets R_i + R_j
  \end{displaymath}
  That is, we can only add one row to another, without doing any
  scaling of that row.  We call this an \textit{addition operation}.

  Demonstrate that two matrices $A$ and $B$ are row equivalent if
  there is a sequence of addition and scaling operations which
  transform $A$ to $B$.  In particular, addition and scaling
  operations can simulate replacement and exchange operations.
\end{tasks}
